% Vietnamese translation for OI Public Library
% Source-File: IOI中国国家候选队论文/2003/邵烜程/邵烜程.ppt
\documentclass[11pt]{article}
\usepackage{oipl-vn}

\title{Bản trình chiếu: Tư tưởng toán học trong thiết kế thuật toán}
\author{Shao Xuancheng}
\date{2003}

\begin{document}
\maketitle

\origtitle{Tư tưởng toán học giúp ta một tay}
\sourcepath{IOI China candidate team papers / 2003 / Shao Xuancheng / Shao Xuancheng.ppt}

\section{Vai trò của toán học}

Bài trình bày mở đầu từ nhận xét: lập trình thi đấu không chỉ là chọn cấu
trúc dữ liệu rồi cài đặt. Nhiều bài có kích thước đủ lớn để tìm kiếm trực
tiếp thất bại, nhưng lại có một quan hệ toán học ẩn giúp biến bài toán thành
một công thức, một hệ phương trình hoặc một cận chặt.

Các slide sắp xếp bốn vai trò của tư tưởng toán học:
\begin{enumerate}
  \item tìm quy luật tổng quát của nghiệm;
  \item xây dựng mô hình toán học để đơn giản hóa đề bài;
  \item dùng phân tích toán học để biến chứng minh thành thuật toán cấu tạo;
  \item dùng kết luận toán học như một cận để tối ưu tìm kiếm.
\end{enumerate}
Những ví dụ như \textit{NOI 2002}, bài cầu, bài đèn tam giác và bài phủ bằng
quân hình chữ L đều nhằm minh họa cùng một điều: trước khi viết chương trình,
cần tìm cấu trúc làm giảm số lựa chọn thực sự.

\section{Ví dụ 1: phân tích số để tích lớn nhất}

Bài toán: cho số nguyên dương \(N\), phân tích \(N\) thành tổng các số tự
nhiên đôi một khác nhau sao cho tích của chúng lớn nhất. Slide dùng ví dụ
\(N=55\) để dẫn tới quy luật.

Nếu viết
\[
  N=a_1+a_2+\cdots+a_m,\qquad
  0<a_1<a_2<\cdots<a_m,
\]
thì phương án tốt phải có các số hạng càng gần nhau càng tốt. Số \(1\) không
nên xuất hiện, vì nó không làm tăng tích và có thể gộp vào số khác. Vì vậy ta
bắt đầu từ dãy
\[
  2,3,\ldots,k
\]
với \(k\) lớn nhất sao cho tổng không vượt quá \(N\). Phần dư còn lại được
phân bổ vào các số hạng lớn ở cuối dãy để giữ dãy gần liên tiếp.

Phần chứng minh dùng lập luận trao đổi: nếu trong phương án có hai số hạng
cách nhau ít nhất \(2\), chuyển một đơn vị từ số lớn sang số nhỏ sẽ làm tích
tăng hoặc không giảm. Sau khi quy luật được chứng minh, thuật toán chỉ còn là
xây dãy và nhân số lớn, thay vì quy hoạch động \(O(N^2)\). Slide ghi độ phức
tạp sau rút gọn là \(O(N)\), phần tốn kém chủ yếu nằm ở xử lý số nguyên lớn.

\section{Ví dụ 2: tháp đèn tam giác}

Ví dụ thứ hai cho một tháp đèn hình tam giác. Trạng thái của đèn ở hàng trên
được quyết định bởi hai đèn phía dưới. Các ràng buộc trong slide có dạng
\[
  3\le N\le1000,\qquad 0<N\le50,\qquad 1\le j\le i<N,
\]
và nếu không có cấu hình thỏa mãn thì in \texttt{NO ANSWER!}.

Điểm chuyển hóa là coi tắt/sáng là \(0/1\), còn quy tắc đèn là phép cộng
modulo \(2\):
\[
  \text{light}[i,j]=
  \text{light}[i+1,j]\oplus\text{light}[i+1,j+1].
\]
Khi đó mọi đèn trong tháp đều là một biểu thức tuyến tính theo các biến của
hàng đáy
\[
  x_1,x_2,\ldots,x_N.
\]
Mỗi đèn đã biết tạo thành một phương trình tuyến tính trên trường
\(\mathbb F_2\). Bài toán đếm cấu hình hàng đáy trở thành bài toán khử Gauss:
nếu hệ vô nghiệm thì in \texttt{NO ANSWER!}; nếu hạng của hệ là \(r\), số
cấu hình là
\[
  2^{N-r}.
\]

Điều đáng học không phải chỉ là kỹ thuật khử Gauss, mà là bước mô hình hóa.
Một bài tìm kiếm \(2^N\) cấu hình được đổi thành một hệ tuyến tính vì ta nhận
ra phép \(\oplus\) nằm sau quy tắc sáng/tắt.

\section{Ví dụ 3: bài toán khóa}

Ví dụ thứ ba có \(M\) nhân viên và yêu cầu thiết kế thẻ sao cho bất kỳ \(N\)
người nào cũng mở được khóa, nhưng bất kỳ \(N-1\) người nào thì không. Mục
tiêu là tối thiểu hóa số loại đặc trưng trên khóa. Slide nhấn mạnh các giới
hạn nhỏ như \(3\le M\le7\), \(1\le N\le4\), \(N\le M\), nhưng lời giải là một
cấu trúc tổ hợp tổng quát.

Xét một nhóm \(N-1\) người. Nhóm này phải thiếu ít nhất một đặc trưng. Nếu
thêm bất kỳ người nào trong \(M-N+1\) người còn lại thì khóa phải mở được,
vì vậy những người còn lại phải cùng sở hữu đặc trưng mà nhóm \(N-1\) người
đang thiếu. Hai nhóm \(N-1\) người khác nhau không thể thiếu cùng một đặc
trưng, nếu không một nhóm đủ \(N\) người vẫn có thể thiếu đặc trưng đó.

Từ chứng minh cận dưới, ta nhận ngay cách dựng:
\begin{enumerate}
  \item xét mọi tập \(X\) gồm \(M-N+1\) người;
  \item tạo một đặc trưng riêng cho \(X\);
  \item ghi đặc trưng đó lên thẻ của tất cả người trong \(X\).
\end{enumerate}
Số đặc trưng cần dùng là
\[
  \binom{M}{N-1}=\binom{M}{M-N+1}.
\]
Đây là ví dụ điển hình cho việc chứng minh và thuật toán không tách rời nhau:
câu ``nhóm còn lại cùng sở hữu một đặc trưng'' chính là thuật toán xây dựng.

\section{Ví dụ 4: phủ bằng quân ba ô hình chữ L}

Ví dụ cuối cùng là một bài tìm kiếm. Cho bảng \(m\times n\), cần đặt nhiều
quân ba ô hình chữ L nhất, hay tương đương làm số ô trống nhỏ nhất. Nếu quay
lui thô, mỗi ô chưa xử lý tạo ra nhiều nhánh: đặt quân theo một hướng hợp lệ
hoặc bỏ trống ô đó.

Tư tưởng toán học ở đây là tính trước một cận đủ chặt cho số ô trống tối ưu.
Các trường hợp diện tích và biên nhỏ cho biết số quân tối đa có thể đặt hoặc
ít nhất một giới hạn không thể vượt qua. Khi đang quay lui, nếu số ô đã bỏ
trống vượt quá cận cho phép, nhánh hiện tại bị loại ngay.

Quy trình trong slide có thể đọc thành:
\begin{enumerate}
  \item duyệt ô theo thứ tự cố định;
  \item thử mọi cách đặt quân hình chữ L hợp lệ liên quan tới ô hiện tại;
  \item cập nhật bảng và số quân đã đặt;
  \item thử nhánh để ô hiện tại trống;
  \item dùng cận toán học để cắt nhánh sớm.
\end{enumerate}
Ở ví dụ này, toán học không thay thế hoàn toàn tìm kiếm; nó làm cho tìm kiếm
có thể chạy được bằng cách loại bỏ các nhánh chắc chắn không tối ưu.

\section{Tổng kết}

Bốn ví dụ tương ứng với bốn kiểu hỗ trợ khác nhau. Có lúc toán học cho công
thức trực tiếp, có lúc cho mô hình đại số, có lúc tạo ra thuật toán trong quá
trình chứng minh, và có lúc chỉ cung cấp cận để quay lui hiệu quả hơn. Điểm
chung là ta phải nhìn bài toán trước khi cài đặt: nếu chỉ bắt đầu từ chương
trình, cấu trúc quan trọng rất dễ bị che khuất.

\end{document}
